\chapter{Complete Implementation of SLH-DSA (FIPS 205)}
\chaptermark{SLH-DSA (FIPS 205)}

\section{Learning objectives}

ML-DSA in the previous chapter rested on lattice problems. FIPS 205, \textbf{SLH-DSA} (from the submission SPHINCS\textsuperscript{+})~\cite{fips205,sphincsplus2019}, rests on nothing but hash functions. There is no number theory, no lattice, no algebraic structure an attacker could exploit -- if the hash function is secure, so is the signature. This makes it the most conservative post-quantum signature, chosen precisely as a hedge in case a structural weakness is ever found in lattices.

After finishing this chapter, you should be able to:

\begin{enumerate}
  \item explain why a secure hash function alone is enough to build a signature;
  \item build up from a one-time signature (WOTS\textsuperscript{+}) to a many-time signature via a Merkle tree;
  \item explain the few-time signature FORS and how it signs the message digest;
  \item see how a hypertree of Merkle trees makes the scheme stateless and practical;
  \item follow the SLH-DSA key generation, signing, and verification flow.
\end{enumerate}

\section{Review: why build a signature from hashes alone?}

Every scheme so far hid a secret inside a structured hard problem -- factoring, lattices, codes. Hash-based signatures take a completely different stance. Their only assumption is that a cryptographic hash $H$ is \emph{one-way} (hard to invert) and \emph{collision-resistant}. Even a quantum computer only speeds up these attacks quadratically (Grover's algorithm), which is countered simply by doubling the output length. There is no Shor-style break, because there is no hidden algebraic structure to exploit.

The whole chapter is a ladder. Each rung fixes a limitation of the one below:

\begin{enumerate}
  \item \textbf{WOTS\textsuperscript{+}} -- signs one message, once.
  \item \textbf{Merkle tree (XMSS)} -- combines many one-time keys under a single public root.
  \item \textbf{FORS} -- a few-time signature used to sign the actual message digest.
  \item \textbf{Hypertree} -- a tree of Merkle trees, which makes the whole thing stateless.
\end{enumerate}

\section{Interlude: what a secure hash function guarantees}

Earlier chapters used hash functions such as SHAKE only as \emph{expanders} -- deterministic ways to stretch a short seed into a matrix or noise (Chapters~\ref{ch:cbd}--\ref{ch:fips203}). This chapter uses them as the \emph{sole source of security}, so it is worth stating precisely what a cryptographic hash $H:\{0,1\}^* \to \{0,1\}^n$ is assumed to provide:

\begin{itemize}
  \item \textbf{Preimage resistance} (one-wayness): given an output $y$, it is hard to find any $x$ with $H(x)=y$.
  \item \textbf{Second-preimage resistance}: given $x$, it is hard to find a different $x'$ with $H(x')=H(x)$.
  \item \textbf{Collision resistance}: it is hard to find \emph{any} pair $x \neq x'$ with $H(x)=H(x')$.
\end{itemize}

That is the entire assumption. There is no group, no lattice, no number-theoretic structure for an attacker to exploit -- which is exactly why hash-based signatures are considered the most conservative option. It also explains the parameter sizes: by Grover's algorithm (Chapter~\ref{ch:quantum}), a quantum search cuts preimage resistance from $n$ bits to about $n/2$, so SLH-DSA simply uses a large enough output $n$ (16, 24, or 32 bytes) to keep a comfortable margin. The threat that shatters RSA leaves hashing merely needing bigger parameters.

\section{Phase 1: the one-time core (WOTS\textsuperscript{+})}

The atomic building block is a \emph{one-time} signature~\cite{lamport1979}: a key pair that may safely sign exactly one message. WOTS\textsuperscript{+} (Winternitz) works by repeated hashing along chains.

To sign a value between $0$ and $w-1$, start from a secret value $sk$ and hash it forward that many times; the public value is the end of the chain, $pk = H^{w-1}(sk)$. The signature is the intermediate point $H^{m}(sk)$. A verifier hashes the signature the remaining $w-1-m$ times and checks it reaches $pk$.

\begin{figure}[H]
\centering
\begin{tikzpicture}[node distance=10mm, every node/.style={font=\scriptsize}]
  \node[flownode] (sk) {$sk$};
  \node[flownode, right=of sk] (h1) {$H(sk)$};
  \node[flownode, right=of h1] (h2) {$H^2(sk)$};
  \node[right=6mm of h2] (dots) {$\cdots$};
  \node[keynode, right=6mm of dots] (pk) {$pk=H^{w-1}(sk)$};
  \draw[flowarrow] (sk) -- (h1);
  \draw[flowarrow] (h1) -- (h2);
  \draw[flowarrow] (h2) -- (dots);
  \draw[flowarrow] (dots) -- (pk);
  \draw[decorate,decoration={brace,amplitude=4pt},pqcorange]
        (h2.north west) -- node[above=3pt]{signature $=H^{m}(sk)$} (h2.north east);
\end{tikzpicture}
\caption{One WOTS\textsuperscript{+} hash chain of length $w$. The secret is the start, the public value is the end, and a signature reveals an intermediate point. A verifier hashes forward to check it reaches the public end. A full message is split into digits, each signed by its own chain; a checksum over the digits prevents an attacker from ``advancing'' any chain to forge a larger value.}
\label{fig:wots-chain}
\end{figure}

A full message digest is chopped into several base-$w$ digits, each signed by its own chain, plus a short \emph{checksum} of those digits (also signed by chains). The checksum is the trick that stops forgery: to change a digit an attacker would have to move at least one chain \emph{backward}, which means inverting $H$. This one-time key pair is safe for a single signature only; signing twice would reveal enough chain points to forge.

\section{Phase 2: from one-time to many-time (the Merkle tree)}

A one-time key is nearly useless on its own. The Merkle tree~\cite{merkle1987} fixes this: generate many WOTS\textsuperscript{+} key pairs, hash each public key into a leaf, and build a binary tree of hashes up to a single \emph{root}. That one root is the public key for \emph{all} the leaves at once.

To use leaf $i$, the signer publishes the WOTS\textsuperscript{+} signature \emph{plus} the \emph{authentication path}: the sibling hash at each level from the leaf up to the root. The verifier recomputes the leaf, then folds it up the path and checks the final value equals the known root.

\begin{figure}[H]
\centering
\begin{tikzpicture}[level distance=9mm, every node/.style={font=\scriptsize},
  level 1/.style={sibling distance=34mm},
  level 2/.style={sibling distance=17mm},
  level 3/.style={sibling distance=9mm}]
  \node[keynode] (root) {root $=pk$}
    child {node[flownode] {$n_0$}
      child {node[flownode,fill=pqclight] (a) {}
        child {node[latticeptbold,label=below:{\scriptsize leaf}] (L) {}}
        child {node[flownode,draw=pqcorange] (s0) {}}
      }
      child {node[flownode,draw=pqcorange] (s1) {}
        child {node[flownode] {}}
        child {node[flownode] {}}
      }
    }
    child {node[flownode,draw=pqcorange] (n1) {$n_1$}
      child {node[flownode] {}
        child {node[flownode] {}}
        child {node[flownode] {}}
      }
      child {node[flownode] {}
        child {node[flownode] {}}
        child {node[flownode] {}}
      }
    };
\end{tikzpicture}
\caption{A Merkle tree of height $3$. Each leaf is the hash of one WOTS\textsuperscript{+} public key; internal nodes hash their two children; the single root is the public key. To authenticate the marked leaf, the signature carries the three orange sibling nodes (the authentication path). The verifier recomputes the leaf and hashes up the path, accepting only if it lands on the known root.}
\label{fig:merkle}
\end{figure}

\begin{algorithm}[H]
\caption{Verify a Merkle authentication path}
\label{alg:merkle-path}
\begin{algorithmic}[1]
\Require Leaf value $\ell$, leaf index $i$, authentication path $(a_0,\dots,a_{h-1})$, root $r$
\Ensure \emph{accept} or \emph{reject}
\State $node \gets \ell$
\For{$j = 0,\dots,h-1$}
  \If{bit $j$ of $i$ is $0$}
    \State $node \gets H(node \,\|\, a_j)$ \Comment{current node is the left child}
  \Else
    \State $node \gets H(a_j \,\|\, node)$ \Comment{current node is the right child}
  \EndIf
\EndFor
\State \Return \emph{accept} if $node = r$, else \emph{reject}
\end{algorithmic}
\end{algorithm}

The catch: each leaf is still a \emph{one-time} key, so the signer must never reuse a leaf. That is a \emph{stateful} requirement -- the signer has to remember which leaves are spent. SLH-DSA removes this burden in the phases below.

\section{Phase 3: signing the message with FORS}

The actual message digest is not signed by a single WOTS\textsuperscript{+} leaf but by \textbf{FORS} (Forest Of Random Subsets), a \emph{few-time} signature. FORS is a small forest of $k$ Merkle trees, each of height $a$. The message digest is chopped into $k$ indices; index $j$ selects one leaf in tree $j$, and the signature reveals that leaf's secret value together with its authentication path to the tree's root. The $k$ tree roots are hashed together into the single FORS public key.

FORS tolerates a few reuses (unlike a one-time key), which is exactly what statelessness needs: because leaves are chosen pseudorandomly from the message, the occasional collision is survivable rather than fatal.

\section{Phase 4: the hypertree makes it stateless}

If there were only one Merkle tree, it would have to be astronomically tall to provide enough leaves to never repeat -- impractical to build. SLH-DSA instead uses a \textbf{hypertree}: $d$ layers of Merkle trees stacked so that each tree's root is signed by a WOTS\textsuperscript{+} key in the tree one layer up. Only the very top root is published as the public key.

\begin{figure}[H]
\centering
\begin{tikzpicture}[every node/.style={font=\scriptsize}]
  \node[keynode] (top) at (0,3) {top tree \; root $=PK.\mathrm{root}$};
  \node[flownode] (mid) at (0,1.7) {middle tree (signs the root below with WOTS\textsuperscript{+})};
  \node[flownode] (bot) at (0,0.4) {bottom tree (signs the FORS public key)};
  \node[flownode, draw=pqcorange] (fors) at (0,-0.9) {FORS \;(signs the message digest)};
  \draw[flowarrow] (fors) -- (bot);
  \draw[flowarrow] (bot) -- (mid);
  \draw[flowarrow] (mid) -- (top);
  \node[right=2mm of top, pqcgray] {published};
  \node[right=2mm of fors, pqcgray] {message};
\end{tikzpicture}
\caption{The hypertree (shown with $d=3$ layers). FORS signs the message digest; the bottom Merkle tree signs the FORS public key; each higher layer signs the root of the tree below it; only the top root is the public key. A single ``signature'' is this whole chain of WOTS\textsuperscript{+} signatures and authentication paths from the message all the way to the published root.}
\label{fig:hypertree}
\end{figure}

Statelessness comes from \emph{how the path is chosen}. Instead of a counter, the signer hashes the message (with a random value $R$) to pick which bottom-layer tree and leaf to use, and which FORS leaves to reveal. The hypertree is so enormous that two messages almost never land on the same leaves, and FORS absorbs the rare overlaps. No state is ever recorded.

\section{Key generation, signing, verification}

Everything is derived deterministically from three $n$-byte values: a secret seed \texttt{SK.seed} (generates every WOTS\textsuperscript{+}/FORS secret via a PRF and an address), a secret PRF key \texttt{SK.prf}, and a public seed \texttt{PK.seed} (tweaks every hash so structurally identical nodes differ). The public key is \texttt{(PK.seed, PK.root)}.

\begin{algorithm}[H]
\caption{SLH-DSA.KeyGen (teaching form)}
\label{alg:slhdsa-keygen}
\begin{algorithmic}[1]
\Require Randomness for three $n$-byte seeds
\Ensure Public key $(\texttt{PK.seed},\texttt{PK.root})$; secret key adds $(\texttt{SK.seed},\texttt{SK.prf})$
\State sample \texttt{SK.seed}, \texttt{SK.prf}, \texttt{PK.seed} uniformly
\State $\texttt{PK.root} \gets$ root of the top-layer Merkle tree, derived from \texttt{SK.seed}, \texttt{PK.seed}
\State \Return public $(\texttt{PK.seed},\texttt{PK.root})$, secret $(\texttt{SK.seed},\texttt{SK.prf},\texttt{PK.seed},\texttt{PK.root})$
\end{algorithmic}
\end{algorithm}

\begin{algorithm}[H]
\caption{SLH-DSA.Sign (teaching form)}
\label{alg:slhdsa-sign}
\begin{algorithmic}[1]
\Require Secret key, message $M$
\Ensure Signature $\sigma = (R,\ \sigma_{\mathrm{FORS}},\ \sigma_{\mathrm{HT}})$
\State $R \gets \textsc{PRF}(\texttt{SK.prf}, M)$ \Comment{randomizer}
\State $digest \gets H(R \,\|\, \texttt{PK.seed} \,\|\, \texttt{PK.root} \,\|\, M)$
\State split $digest$ into the FORS message $md$ and indices $(idx_{\mathrm{tree}}, idx_{\mathrm{leaf}})$
\State $\sigma_{\mathrm{FORS}} \gets \textsc{ForsSign}(md, \dots)$;\quad $pk_{\mathrm{FORS}} \gets$ its public key
\State $\sigma_{\mathrm{HT}} \gets \textsc{HyperTreeSign}(pk_{\mathrm{FORS}}, idx_{\mathrm{tree}}, idx_{\mathrm{leaf}}, \dots)$
\State \Return $(R,\ \sigma_{\mathrm{FORS}},\ \sigma_{\mathrm{HT}})$
\end{algorithmic}
\end{algorithm}

\begin{algorithm}[H]
\caption{SLH-DSA.Verify (teaching form)}
\label{alg:slhdsa-verify}
\begin{algorithmic}[1]
\Require Public key $(\texttt{PK.seed},\texttt{PK.root})$, message $M$, signature $(R,\sigma_{\mathrm{FORS}},\sigma_{\mathrm{HT}})$
\Ensure \emph{accept} or \emph{reject}
\State $digest \gets H(R \,\|\, \texttt{PK.seed} \,\|\, \texttt{PK.root} \,\|\, M)$; split into $md,(idx_{\mathrm{tree}},idx_{\mathrm{leaf}})$
\State $pk_{\mathrm{FORS}} \gets \textsc{ForsPkFromSig}(md, \sigma_{\mathrm{FORS}}, \dots)$
\State $root' \gets \textsc{HyperTreeRootFromSig}(pk_{\mathrm{FORS}}, \sigma_{\mathrm{HT}}, idx_{\mathrm{tree}}, idx_{\mathrm{leaf}}, \dots)$
\State \Return \emph{accept} if $root' = \texttt{PK.root}$, else \emph{reject}
\end{algorithmic}
\end{algorithm}

Verification is entirely a chain of hash recomputations: rebuild the FORS public key from the message, then fold it up the hypertree using the supplied authentication paths, and accept only if the final value equals the published root.

\section{Parameter sets}

SLH-DSA comes in two flavours at each security level: \textbf{s} (\emph{small} signatures, slower signing) and \textbf{f} (\emph{fast} signing, larger signatures), instantiated with either SHA2 or SHAKE.

\begin{table}[H]
\centering
\begin{tabular}{lccc}
\toprule
Family & $n$ (bytes) & NIST level & Trade-off \\
\midrule
SLH-DSA-128\{s,f\} & 16 & 1 & smallest keys \\
SLH-DSA-192\{s,f\} & 24 & 3 & balanced \\
SLH-DSA-256\{s,f\} & 32 & 5 & highest security \\
\bottomrule
\end{tabular}
\caption{SLH-DSA parameter families. Within each level, the \textbf{s} variant minimizes signature size and the \textbf{f} variant minimizes signing time.}
\end{table}

Signatures are large -- roughly $8$ to $50$ kilobytes -- which is the price of relying on hashing alone. In return, the security argument is the simplest and most conservative of any post-quantum scheme.

\section{Pitfalls}

\paragraph{Pitfall 1: reusing a one-time (WOTS\textsuperscript{+}) leaf.}
Signing two different messages with the same WOTS\textsuperscript{+} chain reveals enough chain points to forge. The entire tree machinery exists to guarantee each leaf is used at most once; FORS is deliberately \emph{few-time}, not many-time.

\paragraph{Pitfall 2: thinking SLH-DSA needs state.}
Earlier hash-based schemes (XMSS, LMS) are \emph{stateful}: the signer must remember which leaves are spent, and a single reuse is catastrophic. SLH-DSA is \emph{stateless} -- the message hash selects the path -- which is why it was chosen for standardization despite larger signatures.

\paragraph{Pitfall 3: forgetting the public seed's role.}
\texttt{PK.seed} tweaks every hash call. Without it, structurally identical nodes in different trees would hash identically, enabling multi-target attacks. It is public but essential.

\section{Summary}

SLH-DSA is a ladder of hash constructions:

\begin{itemize}
  \item \textbf{WOTS\textsuperscript{+}} signs once, using hash chains and a checksum;
  \item a \textbf{Merkle tree} unites many one-time keys under one root via authentication paths;
  \item \textbf{FORS} is a few-time signature that signs the message digest;
  \item a \textbf{hypertree} stacks trees so only the top root is public, and hashing the message to pick the path makes the whole scheme stateless.
\end{itemize}

Its security depends on nothing but the hash function, making it the conservative backstop among post-quantum signatures. The next chapter returns to lattices for the third signature standard, FIPS 206 (FN-DSA / Falcon), which trades this simplicity for very compact signatures.
